% 第2章：任务建模与经典规划
\chapter{任务建模与经典规划}

\begin{levelbox}
本章介绍任务规划的形式化建模方法（STRIPS/PDDL）和基于前向搜索的经典规划算法。通过案例A-v1和案例B-v1学习如何将实际问题抽象为规划问题。建议学时：6学时。
\end{levelbox}

% =============================================
\section{\caseA{v1}单星观测任务：问题抽象}
% =============================================

\begin{caseboxA}[案例A-v1：单星确定性观测任务]
\textbf{场景设定}：
\begin{itemize}
  \item 1颗遥感卫星，需观测地面5个目标点
  \item 每个目标有观测收益值
  \item 卫星按固定轨道运行，每轨可观测若干目标
  \item 目标：选择观测序列，最大化总收益
\end{itemize}

\textbf{简化假设}（v1版本）：
\begin{itemize}
  \item 忽略时间窗口约束（假设所有目标均可观测）
  \item 忽略存储和能量约束
  \item 确定性环境（观测必定成功）
\end{itemize}

\textbf{预习思考}：
\begin{enumerate}
  \item 如何定义卫星观测问题的状态？
  \item 需要哪些动作？动作的前置条件和效果是什么？
  \item 如何表示目标条件？
\end{enumerate}
\end{caseboxA}

\subsection{问题抽象}

将卫星观测问题抽象为规划问题：

\textbf{状态表示}：
\begin{itemize}[leftmargin=2em]
  \item 每个目标点的观测状态：已观测/未观测
  \item 卫星的当前位置（可离散化为轨道位置编号）
\end{itemize}

\textbf{动作定义}：
\begin{itemize}[leftmargin=2em]
  \item \texttt{observe(target)}：观测目标target
  \item \texttt{move(pos1, pos2)}：卫星从位置pos1移动到pos2
\end{itemize}

\textbf{目标}：所有高优先级目标均被观测

% =============================================
\section{STRIPS/PDDL建模语言}
% =============================================

\subsection{STRIPS表示}

STRIPS（Stanford Research Institute Problem Solver）是最经典的规划问题表示语言。

\begin{definition}[STRIPS规划问题]
一个STRIPS规划问题由四元组 $\langle P, O, I, G \rangle$ 定义：
\begin{itemize}
  \item $P$：命题（谓词）集合
  \item $O$：操作（动作）集合，每个操作 $o$ 包含：
    \begin{itemize}
      \item $\text{Precond}(o)$：前置条件（命题集合）
      \item $\text{Add}(o)$：添加效果（动作后为真的命题）
      \item $\text{Del}(o)$：删除效果（动作后为假的命题）
    \end{itemize}
  \item $I \subseteq P$：初始状态
  \item $G \subseteq P$：目标条件
\end{itemize}
\end{definition}

\subsection{PDDL建模}

PDDL（Planning Domain Definition Language）是STRIPS的标准化扩展，是国际规划竞赛（IPC）的标准语言。

\begin{keypoint}[PDDL的两个文件]
\begin{itemize}
  \item \textbf{领域文件}（Domain File）：定义谓词和动作模式，与具体问题实例无关
  \item \textbf{问题文件}（Problem File）：定义具体的对象、初始状态和目标
\end{itemize}
\end{keypoint}

\subsection{案例A-v1的PDDL建模}

\textbf{领域文件}（satellite-domain.pddl）：

\begin{lstlisting}[language=Lisp,caption=卫星观测领域定义]
(define (domain satellite-observe)
  (:requirements :strips)

  (:predicates
    (satellite ?s)           ; s是卫星
    (target ?t)              ; t是目标
    (position ?p)            ; p是位置
    (at ?s ?p)               ; 卫星s在位置p
    (visible ?t ?p)          ; 目标t在位置p可见
    (observed ?t)            ; 目标t已被观测
  )

  (:action observe
    :parameters (?s ?t ?p)
    :precondition (and (satellite ?s) (target ?t) (position ?p)
                       (at ?s ?p) (visible ?t ?p) (not (observed ?t)))
    :effect (observed ?t)
  )

  (:action move
    :parameters (?s ?from ?to)
    :precondition (and (satellite ?s) (position ?from) (position ?to)
                       (at ?s ?from))
    :effect (and (at ?s ?to) (not (at ?s ?from)))
  )
)
\end{lstlisting}

\textbf{问题文件}（satellite-problem.pddl）：

\begin{lstlisting}[language=Lisp,caption=卫星观测问题实例]
(define (problem observe-5targets)
  (:domain satellite-observe)

  (:objects
    sat1 - satellite
    t1 t2 t3 t4 t5 - target
    p1 p2 p3 p4 p5 - position
  )

  (:init
    (satellite sat1)
    (target t1) (target t2) (target t3) (target t4) (target t5)
    (position p1) (position p2) (position p3) (position p4) (position p5)
    (at sat1 p1)
    (visible t1 p1) (visible t2 p2) (visible t3 p3)
    (visible t4 p4) (visible t5 p5)
  )

  (:goal (and (observed t1) (observed t2) (observed t3)))
)
\end{lstlisting}

% =============================================
\section{前向搜索与启发式规划}
% =============================================

\subsection{前向状态空间搜索}

前向搜索从初始状态出发，逐步应用动作，直到达到目标状态。

\begin{algorithm}[htbp]
\caption{前向状态空间搜索}
\begin{algorithmic}[1]
\Require 规划问题 $\langle P, O, I, G \rangle$
\Ensure 动作序列或失败
\State $\text{Open} \gets \{I\}$，$\text{Closed} \gets \emptyset$
\While{$\text{Open} \neq \emptyset$}
  \State $s \gets$ 从Open中选择状态（按启发式值）
  \If{$G \subseteq s$}
    \State \Return 从$I$到$s$的动作序列
  \EndIf
  \State $\text{Closed} \gets \text{Closed} \cup \{s\}$
  \For{每个可应用的动作 $a \in O$}
    \State $s' \gets \text{apply}(s, a)$
    \If{$s' \notin \text{Closed}$}
      \State 将 $s'$ 加入 Open
    \EndIf
  \EndFor
\EndWhile
\State \Return 失败
\end{algorithmic}
\end{algorithm}

\subsection{规划启发式}

\subsubsection{删除松弛启发式}

\textbf{删除松弛}（Delete Relaxation）：忽略动作的删除效果，只保留添加效果。松弛问题更容易求解，其最优解代价是原问题的下界。

\begin{itemize}[leftmargin=2em]
  \item $h^{\text{add}}$：松弛问题中所有子目标代价之和
  \item $h^{\text{max}}$：松弛问题中代价最大的子目标
  \item $h^{\text{FF}}$：基于规划图的松弛启发式，FF规划器使用
\end{itemize}

\subsubsection{Fast Downward与LAMA}

\textbf{Fast Downward}是现代最成功的规划器之一，采用：
\begin{itemize}[leftmargin=2em]
  \item SAS+表示：多值变量替代命题
  \item 因果图分析：识别变量间的依赖关系
  \item 多种启发式：$h^{\text{FF}}$、因果图启发式、地标启发式
\end{itemize}

\textbf{LAMA}基于地标启发式，多次获得IPC冠军。

% =============================================
\section{\caseB{v1}单平台侦察任务建模}
% =============================================

\begin{caseboxB}[案例B-v1：单无人机侦察任务]
\textbf{场景设定}：
\begin{itemize}
  \item 1架无人机，需侦察3个目标区域
  \item 无人机从基地起飞，完成侦察后返回
  \item 每个目标区域有侦察价值
  \item 目标：规划飞行路线，完成所有侦察任务
\end{itemize}

\textbf{简化假设}（v1版本）：
\begin{itemize}
  \item 忽略航程约束
  \item 忽略敌方威胁
  \item 确定性环境
\end{itemize}
\end{caseboxB}

\subsection{问题建模}

\textbf{状态表示}：
\begin{itemize}[leftmargin=2em]
  \item 无人机位置
  \item 各目标的侦察状态（已侦察/未侦察）
\end{itemize}

\textbf{动作}：
\begin{itemize}[leftmargin=2em]
  \item \texttt{fly(from, to)}：飞行
  \item \texttt{recon(target)}：侦察目标
  \item \texttt{return\_base()}：返回基地
\end{itemize}

\subsection{PDDL建模}

\begin{lstlisting}[language=Lisp,caption=无人机侦察领域定义]
(define (domain uav-recon)
  (:requirements :strips)

  (:predicates
    (uav ?u)
    (location ?l)
    (target ?t)
    (at ?u ?l)
    (target-at ?t ?l)
    (reconned ?t)
    (at-base ?u)
  )

  (:action fly
    :parameters (?u ?from ?to)
    :precondition (and (uav ?u) (location ?from) (location ?to)
                       (at ?u ?from))
    :effect (and (at ?u ?to) (not (at ?u ?from)) (not (at-base ?u)))
  )

  (:action recon
    :parameters (?u ?t ?l)
    :precondition (and (uav ?u) (target ?t) (location ?l)
                       (at ?u ?l) (target-at ?t ?l) (not (reconned ?t)))
    :effect (reconned ?t)
  )

  (:action return-base
    :parameters (?u ?l)
    :precondition (and (uav ?u) (location ?l) (at ?u ?l))
    :effect (at-base ?u)
  )
)
\end{lstlisting}

% =============================================
\section{本章小结与讨论}
% =============================================

\subsection*{本章小结}

\begin{itemize}[leftmargin=2em]
  \item STRIPS/PDDL是任务规划的标准建模语言
  \item 前向状态空间搜索配合启发式函数是经典规划的主流方法
  \item 删除松弛是设计可纳启发式的有效技术
  \item Fast Downward和LAMA是现代高效规划器的代表
\end{itemize}

\begin{discussbox}
\textbf{讨论题}：
\begin{enumerate}
  \item 案例A-v1的PDDL建模忽略了哪些实际约束？这些约束能用基本PDDL表达吗？
  \item 如果目标数量从5个增加到100个，搜索效率会如何变化？有什么改进方法？
  \item 案例B的"返回基地"是否必须作为显式目标？如果不返回会怎样？
  \item 尝试为案例A或案例B设计一个简单的启发式函数，分析其可纳性。
\end{enumerate}

\textbf{编程练习}：
\begin{enumerate}
  \item 使用Fast Downward求解案例A-v1的PDDL问题
  \item 增加目标数量，测试求解时间的变化趋势
  \item 修改目标条件（部分观测即可），比较规划结果的差异
\end{enumerate}
\end{discussbox}
